
\documentclass[11pt,a4paper]{article}
\usepackage[margin=1.9cm]{geometry}
\usepackage{fontspec}
\usepackage{polyglossia}
\setdefaultlanguage{greek}
\setotherlanguage{english}
\setmainfont{DejaVu Serif}[ItalicFont={DejaVu Serif},ItalicFeatures={FakeSlant=0.2},BoldItalicFont={DejaVu Serif Bold},BoldItalicFeatures={FakeSlant=0.2}]
\setmonofont{DejaVu Sans Mono}[Scale=0.88]
\newfontfamily\greekfont{DejaVu Serif}[ItalicFont={DejaVu Serif},ItalicFeatures={FakeSlant=0.2},BoldItalicFont={DejaVu Serif Bold},BoldItalicFeatures={FakeSlant=0.2}]
\newfontfamily\greekfonttt{DejaVu Sans Mono}[Scale=0.88]
\usepackage{amsmath,amssymb,amsthm}
\usepackage{enumitem}
\usepackage{fancyvrb}
\usepackage{fancyhdr}
\usepackage[unicode,hidelinks]{hyperref}
\newcommand{\horrule}[1]{\rule{\linewidth}{#1}}
\newcommand{\eps}{\varepsilon}

\theoremstyle{definition}
\newtheorem{question}{Άσκηση}
\setlist[enumerate]{label=(\alph*),leftmargin=*,itemsep=3pt,topsep=5pt}
\setlength{\parindent}{0pt}
\setlength{\parskip}{5pt}
\setlength{\emergencystretch}{3em}
\pagestyle{fancy}
\fancyhf{}
\fancyfoot[L]{\small ΗΥ280 (2026) - Τσουρακάκης}
\fancyfoot[R]{\small\thepage}
\renewcommand{\headrulewidth}{0pt}
\DefineVerbatimEnvironment{pseudocode}{Verbatim}{fontsize=\small,xleftmargin=1em}


\begin{document}
\begin{center}
\horrule{0.4pt}\\[0.15cm]
{\large\textbf{ΗΥ 280 \textemdash{} Θεωρία Υπολογισμού}}\\[0.1cm]
{Ασκήσεις 5 \textemdash{} Συμβολοσειρές}\\[0.1cm]
{Χαράλαμπος Ε. Τσουρακάκης}\\
\horrule{0.4pt}
\end{center}

\begin{center}\fbox{\begin{minipage}{0.93\linewidth}
\textbf{Οδηγίες για τις ασκήσεις}
\begin{itemize}[leftmargin=*,itemsep=6pt]
\item \textbf{Ασκήσεις 1-4:} Αποτελούν μέρος του παραδοτέου.
\item Στην Άσκηση 1, για τις ερωτήσεις Σωστό/Λάθος δεν χρειάζεται να γράψετε αιτιολόγηση. Θα πρέπει, ωστόσο, να γνωρίζετε την αιτιολόγηση της απάντησής σας.
\item Όπου ζητείται επαγωγή, γράψτε καθαρά βάση, επαγωγική υπόθεση και βήμα.
\item {\bf Αξιοποιήστε τις ώρες γραφείου} αφού πρώτα προσπαθήσετε τις ασκήσεις.
\end{itemize}
\end{minipage}}\end{center}\medskip

\begin{question}[$\bigstar$ Έλεγχος κατανόησης]
Έστω $\Sigma=\{0,1\}$ και $x,y,z\in\Sigma^*$.

\textbf{Σωστό ή Λάθος.}
\begin{enumerate}
\item Αν $|xy|=|xz|$, τότε $y=z$.
\item Αν $xy=xz$, τότε $y=z$.
\item Για κάθε $x,y$, $(xy)^R=x^Ry^R$.
\item Αν $x$ είναι prefix της $y$, τότε $x^R$ είναι suffix της $y^R$.
\item Αν $xy=yx$, τότε $x=y$.
\item Αν $xyz=\eps$, τότε $x=y=z=\eps$.
\end{enumerate}
\end{question}

\begin{question}[$\bigstar$ Η αναδρομική συνάρτηση \emph{hanoi}]
Ορίζουμε
\[
\operatorname{hanoi}(w)=
\begin{cases}
\eps,&w=\eps,\\
\operatorname{hanoi}(x)\,a\,\operatorname{hanoi}(x),&w=ax.
\end{cases}
\]
\begin{enumerate}
\item Υπολογίστε $\operatorname{hanoi}(0)$, $\operatorname{hanoi}(01)$ και $\operatorname{hanoi}(010)$.
\item  Αποδείξτε με επαγωγή ότι
\[
|\operatorname{hanoi}(w)|=2^{|w|}-1
\]
για κάθε $w$.
\end{enumerate}
\end{question}

\begin{question}[$\bigstar\bigstar$ Αναδρομικά ορισμένη γλώσσα]
Ορίζουμε $L\subseteq\{0,1\}^*$ ως τη μικρότερη γλώσσα που ικανοποιεί:
\begin{itemize}
\item $\eps\in L$,
\item αν $x,y\in L$, τότε $0x1y0\in L$.
\end{itemize}
\begin{enumerate}
\item Δείξτε ότι $010\in L$, $001010\in L$ και $00101010010\in L$.
\item Αποδείξτε με δομική επαγωγή ότι κάθε $w\in L$ ικανοποιεί
\[
\#_0(w)=2\#_1(w).
\]
Μπορείτε να χρησιμοποιήσετε $\#_a(xy)=\#_a(x)+\#_a(y)$.
\item Δώστε συμβολοσειρά με διπλάσια $0$ από $1$ που δεν ανήκει στη $L$.
\end{enumerate}
\end{question}

\begin{question}[$\bigstar\bigstar$ Η συνάρτηση καταμέτρησης]
Για σύμβολο $a\in\Sigma$ και $w\in\Sigma^*$, έστω $\#(a,w)$ το πλήθος εμφανίσεων του $a$ στη $w$.
\begin{enumerate}
\item Δώστε αναδρομικό ορισμό της
\[
\#:\Sigma\times\Sigma^*\to\mathbb N.
\]
\item Αποδείξτε ότι για κάθε $a\in\Sigma$ και $x,y\in\Sigma^*$,
\[
\#(a,xy)=\#(a,x)+\#(a,y).
\]
\item Αποδείξτε ότι για κάθε $w\in\Sigma^*$,
\[
|w|=\sum_{a\in\Sigma}\#(a,w).
\]
\end{enumerate}
\end{question}

\end{document}
